\documentclass[11pt]{article}

\usepackage[margin=1in]{geometry}
\usepackage{amsmath,amssymb,mathtools}
\usepackage{booktabs}
\usepackage{enumitem}
\usepackage{hyperref}

\hypersetup{
  colorlinks=true,
  linkcolor=blue,
  urlcolor=blue,
  citecolor=blue
}

\title{OpenMC Delayed-Neutron Import:\\
Consistency Checks and User Reporting}
\author{OpenSn Developer Note}
\date{\today}

\begin{document}
\maketitle

\section{Purpose}

This note defines the consistency checks that should be performed when importing
delayed-neutron multigroup data from an OpenMC MGXS library into OpenSn, and it
defines what information should be reported to the user.

The goal is not to force noisy Monte Carlo data to satisfy deterministic
identities exactly. The goal is to:
\begin{enumerate}[label=\arabic*.]
  \item detect important inconsistencies between steady-state and transient data,
  \item quantify those inconsistencies,
  \item make the approximations used by OpenSn explicit, and
  \item warn the user when a null transient should not be expected.
\end{enumerate}

\section{Background}

OpenSn steady-state transport currently assumes a rank-1 fission operator of the form
\begin{equation}
  F_{g \leftarrow g'}^{\mathrm{SS}} \approx \chi_g \, \nu \Sigma_{f,g'}^{\mathrm{tot}},
\end{equation}
while the transient treatment with delayed neutrons uses prompt production plus
delayed precursor emission:
\begin{align}
  F_{g \leftarrow g'}^{\mathrm{prompt}} &\approx \chi^{p}_g \, \nu \Sigma_{f,g'}^{p}, \\
  \text{delayed source} &\propto \sum_j \chi^{d,j}_g \, \lambda_j C_j.
\end{align}

When OpenMC data are imported, the following issues may prevent a clean null transient:
\begin{enumerate}[label=\arabic*.]
  \item \textbf{Closure may fail due to noise:}
  \begin{equation}
    \nu \Sigma_f^{\mathrm{tot}} \neq \nu \Sigma_f^p + \nu \Sigma_f^d.
  \end{equation}

  \item \textbf{Steady-state chi may not be recoverable exactly} from prompt and delayed
  spectra unless delayed fractions factor cleanly from energy.

  \item \textbf{The true fission operator may not be rank-1.}
  OpenMC may provide a full fission matrix and a rank-1 approximation; OpenSn uses the
  latter representation.
\end{enumerate}

\section{Notation}

For energy group \(g\) and precursor family \(j\), define:
\begin{align}
  \nu \Sigma_{f,g}^{\mathrm{tot}} &: \text{ total neutron production cross section}, \\
  \nu \Sigma_{f,g}^{p} &: \text{ prompt neutron production cross section}, \\
  \nu \Sigma_{f,g}^{d,j} &: \text{ delayed production cross section for family } j, \\
  \nu \Sigma_{f,g}^{d} &= \sum_j \nu \Sigma_{f,g}^{d,j}, \\
  \chi_g &: \text{ imported steady-state fission spectrum}, \\
  \chi^p_g &: \text{ imported prompt fission spectrum}, \\
  \chi^{d,j}_g &: \text{ delayed emission spectrum for family } j, \\
  \beta_g &= \frac{\nu \Sigma_{f,g}^{d}}{\nu \Sigma_{f,g}^{\mathrm{tot}}}, \\
  \beta_{j,g} &= \frac{\nu \Sigma_{f,g}^{d,j}}{\nu \Sigma_{f,g}^{\mathrm{tot}}}, \\
  y_{j,g} &= \frac{\nu \Sigma_{f,g}^{d,j}}{\nu \Sigma_{f,g}^{d}}.
\end{align}

\section{Checks To Perform}

\subsection{Basic validity checks}

These are hard checks. Failure should abort import.
\begin{enumerate}[label=\arabic*.]
  \item All required delayed-neutron datasets are present when delayed groups are declared.
  \item All \(\nu \Sigma_f\), spectra, and decay constants are non-negative.
  \item Each spectrum is either zero everywhere or has positive sum.
  \item Dimensions match the declared number of groups and precursor families.
\end{enumerate}

\subsection{Nu-fission closure checks}

For each group \(g\), compute the residual
\begin{equation}
  r^{\nu}_g = \nu \Sigma_{f,g}^{\mathrm{tot}} -
  \nu \Sigma_{f,g}^{p} -
  \sum_j \nu \Sigma_{f,g}^{d,j}.
\end{equation}

Also compute a relative residual
\begin{equation}
  \epsilon^{\nu}_g =
  \frac{\left|r^{\nu}_g\right|}
       {\max\!\left(\left|\nu \Sigma_{f,g}^{\mathrm{tot}}\right|, \epsilon_{\mathrm{floor}}\right)}.
\end{equation}

In addition, compute global closure:
\begin{equation}
  R^{\nu} =
  \sum_g \nu \Sigma_{f,g}^{\mathrm{tot}} -
  \sum_g \nu \Sigma_{f,g}^{p} -
  \sum_{g,j} \nu \Sigma_{f,g}^{d,j}.
\end{equation}

This is the most important consistency check for transient behavior.

\subsection{Spectrum normalization checks}

For each available spectrum, report its normalization residual:
\begin{align}
  r^{\chi} &= \left|\sum_g \chi_g - 1\right|, \\
  r^{\chi^p} &= \left|\sum_g \chi^p_g - 1\right|, \\
  r^{\chi^{d,j}} &= \left|\sum_g \chi^{d,j}_g - 1\right|.
\end{align}

If a spectrum is not normalized but has positive sum, OpenSn may renormalize it on import,
but the renormalization should be reported explicitly.

\subsection{Energy dependence of delayed fraction}

To determine whether the steady-state spectrum can be approximated from prompt and delayed
data in a simple way, measure the variation of
\begin{equation}
  \beta_g = \frac{\nu \Sigma_{f,g}^{d}}{\nu \Sigma_{f,g}^{\mathrm{tot}}}.
\end{equation}

Recommended reported metrics:
\begin{align}
  \beta_{\min} &= \min_g \beta_g, \\
  \beta_{\max} &= \max_g \beta_g, \\
  \overline{\beta} &= \frac{\sum_g \nu \Sigma_{f,g}^{\mathrm{tot}} \beta_g}
                           {\sum_g \nu \Sigma_{f,g}^{\mathrm{tot}}}, \\
  \Delta \beta &= \beta_{\max} - \beta_{\min}.
\end{align}

If \(\Delta \beta\) is small, then the delayed fraction is nearly energy independent.
If it is not small, then the usual factorized interpretation of \(\chi\) becomes weak.

\subsection{Energy dependence of precursor-family yields}

Measure whether family fractions are approximately independent of incident energy:
\begin{equation}
  y_{j,g} = \frac{\nu \Sigma_{f,g}^{d,j}}{\nu \Sigma_{f,g}^{d}}.
\end{equation}

For each family \(j\), report
\begin{align}
  y_{j,\min} &= \min_g y_{j,g}, \\
  y_{j,\max} &= \max_g y_{j,g}, \\
  \Delta y_j &= y_{j,\max} - y_{j,\min}.
\end{align}

Large variation in \(y_{j,g}\) indicates that even the delayed-family decomposition does
not factor simply with energy.

\subsection{Approximate steady-state spectrum reconstruction}

If \(\beta_g\) and \(y_{j,g}\) are approximately energy independent, define the weighted
average family fractions
\begin{equation}
  \overline{y}_j =
  \frac{\sum_g \nu \Sigma_{f,g}^{d,j}}
       {\sum_g \nu \Sigma_{f,g}^{d}},
\end{equation}
and an effective delayed spectrum
\begin{equation}
  \chi_g^{d,\mathrm{eff}} = \sum_j \overline{y}_j \chi_g^{d,j}.
\end{equation}

Then reconstruct an approximate steady-state spectrum:
\begin{equation}
  \chi_g^{\mathrm{recon}} =
  (1-\overline{\beta}) \chi_g^p + \overline{\beta}\chi_g^{d,\mathrm{eff}}.
\end{equation}

Compare this against the imported steady-state \(\chi_g\) using
\begin{align}
  \delta^{\chi}_{\infty} &= \max_g \left|\chi_g - \chi_g^{\mathrm{recon}}\right|, \\
  \delta^{\chi}_{1} &= \sum_g \left|\chi_g - \chi_g^{\mathrm{recon}}\right|.
\end{align}

This is not an exact identity in the general case. It is a diagnostic showing whether the
steady-state spectrum is compatible with the prompt/delayed decomposition under the
factorization assumptions used by OpenSn.

\subsection{Prompt-spectrum fallback check}

If OpenMC delayed data are present but \(\chi^p\) is absent and OpenSn falls back to
\(\chi^p \leftarrow \chi\), then this should always generate a warning:

\begin{quote}
  Prompt fission spectrum was not provided by the OpenMC library. OpenSn is using the
  steady-state spectrum as a prompt-spectrum surrogate. This is an approximation and may
  prevent a null transient.
\end{quote}

\subsection{Rank-1 fission-operator check}

If the OpenMC library contains a full fission matrix \(F^{\mathrm{full}}\) in addition to the
rank-1 approximation \(F^{\mathrm{rank1}}\), compare them.

Recommended metrics:
\begin{align}
  \delta^F_{\infty} &= \max_{g,g'} \left|F^{\mathrm{full}}_{g \leftarrow g'} -
                                      F^{\mathrm{rank1}}_{g \leftarrow g'}\right|, \\
  \delta^F_{1,\mathrm{rel}} &=
  \frac{\sum_{g,g'} \left|F^{\mathrm{full}}_{g \leftarrow g'} -
                        F^{\mathrm{rank1}}_{g \leftarrow g'}\right|}
       {\max\!\left(\sum_{g,g'} \left|F^{\mathrm{full}}_{g \leftarrow g'}\right|,
                    \epsilon_{\mathrm{floor}}\right)}.
\end{align}

If only rank-1 data are available, report that fact. If full-matrix data are available and
the mismatch is large, warn that steady-state behavior is already using an approximation
before transient delayed-neutron effects are considered.

\section{What To Report To The User}

For each imported material with delayed-neutron data, print a compact summary such as:

\begin{verbatim}
OpenMC delayed-neutron consistency summary for material "fuel":
  groups = 6, precursor families = 8
  nu-sigma-f closure:
    max abs residual       = 2.1e-5
    max rel residual       = 4.6e-4
    global abs residual    = 7.8e-6
  spectrum normalization:
    chi sum                = 1.000000
    chi_prompt sum         = 0.999991
    max chi_delayed sum err= 3.2e-5
  delayed fraction:
    beta_min, beta_max     = 6.12e-3, 6.58e-3
    beta range             = 4.6e-4
  delayed family variation:
    max family fraction range = 1.1e-3
  chi consistency:
    ||chi - chi_recon||_inf = 8.3e-4
    ||chi - chi_recon||_1   = 2.5e-3
  fission operator:
    using rank-1 approximation from OpenMC
    relative full-vs-rank1 mismatch = 1.7e-2
\end{verbatim}

\subsection{Required warning messages}

Warnings should be issued when any of the following occur:
\begin{enumerate}[label=\arabic*.]
  \item \(\nu \Sigma_f^{\mathrm{tot}}\) does not close well against prompt plus delayed parts.
  \item \(\chi^p\) is missing and \(\chi\) is used as a surrogate.
  \item \(\beta_g\) varies appreciably with energy.
  \item precursor-family fractions vary appreciably with energy.
  \item reconstructed \(\chi^{\mathrm{recon}}\) differs appreciably from imported \(\chi\).
  \item the full fission matrix differs appreciably from the rank-1 approximation.
\end{enumerate}

A top-level warning should state:

\begin{quote}
  Imported OpenMC delayed-neutron data are not fully consistent with the steady-state
  fission data used by OpenSn. A null transient should not be expected.
\end{quote}

\section{Suggested Severity Levels}

The exact thresholds should remain configurable, but the reporting logic should follow:

\begin{center}
\begin{tabular}{@{}lll@{}}
\toprule
Condition & Severity & Action \\
\midrule
Missing required delayed dataset & Error & Abort import \\
Negative XS, decay constant, or spectrum entry & Error & Abort import \\
Bad dimensions & Error & Abort import \\
Small closure / normalization mismatch & Info & Report only \\
Moderate closure or chi mismatch & Warning & Report prominently \\
Large closure or rank-1 mismatch & Warning & State null transient not expected \\
\bottomrule
\end{tabular}
\end{center}

\section{Implementation Guidance}

The import routine should keep three categories distinct:
\begin{enumerate}[label=\arabic*.]
  \item \textbf{Hard validity checks:} data are unusable; fail immediately.
  \item \textbf{Soft consistency checks:} data are usable but internally inconsistent; warn.
  \item \textbf{Approximation disclosures:} data are valid, but OpenSn is using a reduced model;
  disclose the approximation even if numerical mismatch is small.
\end{enumerate}

The implementation should not silently ``fix'' Monte Carlo noise by enforcing exact closure,
because doing so hides the actual data quality from the user. If renormalization is applied to
spectra, it should be reported.

\section{Bottom Line}

The delayed-neutron import should be accepted only after it does all of the following:
\begin{enumerate}[label=\arabic*.]
  \item checks prompt-plus-delayed closure against steady-state \(\nu \Sigma_f\),
  \item checks whether the prompt/delayed decomposition is compatible with the imported
  steady-state \(\chi\),
  \item warns when \(\beta\) or delayed-family fractions are energy dependent,
  \item warns when prompt-spectrum fallback is used,
  \item warns when the OpenMC full fission operator is not well represented by the rank-1
  form consumed by OpenSn, and
  \item reports all of the above clearly to the user at import time.
\end{enumerate}

\end{document}
